\documentclass[11pt]{article}
\usepackage[a4paper,margin=2.2cm]{geometry}
\usepackage[T1]{fontenc}
\usepackage{textcomp}
\usepackage[spanish]{babel}
\usepackage{amsmath,amssymb}
\usepackage{enumitem}
\usepackage{tikz}
\usetikzlibrary{calc,arrows.meta,patterns,decorations.pathmorphing}
\usepackage{xcolor}
\usepackage{multicol}
\usepackage{parskip}
\setlength{\parskip}{6pt}
\usepackage{lmodern}
\renewcommand{\familydefault}{\sfdefault}

\definecolor{unalblue}{RGB}{31,73,124}

\setlist[enumerate,1]{itemsep=10pt,topsep=8pt,leftmargin=2.2em,label=\protect\fbox{\arabic*}}
\newlist{opciones}{enumerate}{1}
\setlist[opciones]{label=(\alph*),itemsep=8pt,topsep=6pt,leftmargin=2em,font=\normalfont}

\newcommand{\ptsbox}[1]{\fboxsep=3pt\fbox{\footnotesize #1}\hspace{4pt}}
\newcommand{\borradorbox}[1][3cm]{%
  \noindent\fbox{\begin{minipage}[t][#1][t]{0.97\linewidth}
  {\scriptsize\itshape Espacio borrador, nada escrito aquí será calificado.}
  \end{minipage}}
}

\newcommand{\vect}[2]{\begin{pmatrix}#1\\#2\end{pmatrix}}
\newcommand{\vectiii}[3]{\begin{pmatrix}#1\\#2\\#3\end{pmatrix}}

\newcommand{\examheader}{%
  \noindent
  \begin{minipage}[t]{0.6\linewidth}
    {\small\bfseries Geometría Vectorial y Analítica --- Supletorio Parcial 2}\\
    {\small 16 de Marzo de 2020}
  \end{minipage}%
  \hfill
  \begin{minipage}[t]{0.37\linewidth}
    \raggedleft
    \fbox{\begin{minipage}{0.95\linewidth}\centering\scriptsize Escuela de Matemáticas. Universidad Nacional de Colombia, Sede Medellín.\end{minipage}}
  \end{minipage}\\[4pt]
  \rule{\linewidth}{0.6pt}
}
\newcommand{\examfooter}{%
  \vfill
  \rule{\linewidth}{0.4pt}\\[1pt]
  {\scriptsize\textit{Transcripción digital --- MathUNAL}\hfill\textcolor{unalblue}{\textbf{\thepage}}}
}
\pagestyle{empty}

\begin{document}
\examheader

\begin{center}
{\bfseries Supletorio del Segundo Parcial de Geometría Vectorial}\\
16 de Marzo de 2020

\vspace{6pt}
\textbf{Puntaje. Sólo para uso Oficial}

\vspace{4pt}
\begin{tabular}{|c|c|c|c|c|c|c|}
\hline
1-4 (/30) & 5 (/20) & 6 (/20) & 7 (/20) & 8 (/10) & Total & Nota\\
\hline
 & & & & & & \\
\hline
\end{tabular}
\end{center}

\textbf{Instrucciones:} La duración del examen es de \textbf{1 hora y 50 minutos}. El examen consta de ocho preguntas en dos hojas impresas por ambos lados. Antes de empezar verifique que su examen esté completo y que sus datos sean correctos. No desengrape su examen ni añada otras grapas. Utilice los espacios asignados para resolver cada pregunta. \textbf{No} está permitido: hacer preguntas, usar calculadoras, sacar hojas en blanco, ni tomar apuntes en hojas diferentes a las del examen. No se permite que ningún celular se encuentre a la vista.

\vspace{10pt}
En las preguntas 1 a 4 responda en el espacio en blanco. \textbf{En estas preguntas la calificación tendrá en cuenta sólo la respuesta.}

\begin{enumerate}

\item \ptsbox{7\%} Sean $\mathcal{L}_1$ y $\mathcal{L}_2$ las líneas rectas con ecuaciones $\mathcal{L}_1: \vect{x}{y} = \vect{-1}{0} + t\vect{-1}{1}$ y $\mathcal{L}_2: x-y+3=0$. Entonces el ángulo entre $\mathcal{L}_1$ y $\mathcal{L}_2$ es

\vspace{50pt}
\underline{\hspace{6cm}}.

\item \ptsbox{7\%} Si $\mathcal{L}$ es la línea recta con ecuación $x-2y+5=0$ entonces la línea recta que pasa por $\vect{2}{-1}$, y es perpendicular a $\mathcal{L}$, tiene por ecuación general:

\vspace{60pt}
\underline{\hspace{6cm}}.

\item \ptsbox{8\%} Si $T$ es la transformación lineal definida por $T\vect{x}{y} = \vect{-6x+24y}{3x-12y}$ entonces el núcleo de $T$ es la línea recta con ecuación general

\vspace{80pt}
\underline{\hspace{6cm}}.

\item \ptsbox{8\%} Si $T$ es una transformación lineal tal que $T\vect{3}{-4} = \vect{1}{1}$ entonces $T\vect{-15}{20} = $

\vspace{80pt}
\underline{\hspace{6cm}}.

\end{enumerate}

\vspace{10pt}
\noindent\textbf{PREGUNTAS CON PROCEDIMIENTO}

\noindent En las preguntas 5 a 8 responda mostrando \emph{detalladamente} el procedimiento utilizado.

\begin{enumerate}
\setcounter{enumi}{4}

\item \ptsbox{20\%} Sean $S$ y $T$ las transformaciones lineales dadas por $S\vect{x}{y} = \vect{-3x+11y}{-x+4y}$ y $T\vect{x}{y} = \vect{-8x+3y}{5x-2y}$. Halle $(S\circ T)^{-1}\vect{x}{y}$ para todo $\vect{x}{y}$ en el plano.

\vspace{4pt}
{\scriptsize\itshape Emplee este espacio para realizar operaciones.}

\vspace{160pt}

\item \ptsbox{20\%} Considere la recta $\mathcal{L}$ que pasa por los puntos $A = \vect{-3}{5}$ y $B = \vect{-7}{9}$. Halle una ecuación general de la mediatriz del segmento de recta $\overline{AB}$.

\vspace{4pt}
{\scriptsize\itshape Emplee este espacio para realizar operaciones.}

\vspace{160pt}

\newpage
\examheader

\item \ptsbox{20\%} Sea $T$ la transformación lineal que satisface $T\vect{3}{4} = \vect{10}{-1}$ y $T\vect{1}{-5} = \vect{-3}{6}$.

\begin{opciones}
\item \ptsbox{10\%} Halle la matriz de $T$, $m(T)$.

\vspace{4pt}
{\scriptsize\itshape Emplee este espacio para realizar operaciones.}

\vspace{160pt}

\item \ptsbox{10\%} Si $\mathcal{L}$ es la recta con ecuación general $2x-3y+4=0$, halle una ecuación general de la recta $T(\mathcal{L})$.

\vspace{4pt}
{\scriptsize\itshape Emplee este espacio para realizar operaciones.}

\vspace{160pt}
\end{opciones}

\item \ptsbox{10\%} Sea $T$ una transformación lineal de plano, y sean $U$ y $V$ dos vectores linealmente independientes tales que $T(U) = -U$ y $T(V) = -V$. Demuestre que $T(X) = -X$ cualquiera sea $X\in\mathbb{R}^2$.

\vspace{4pt}
{\scriptsize\itshape Emplee este espacio para realizar operaciones.}

\vspace{160pt}

\end{enumerate}

\examfooter
\end{document}
